\section{Discussion, External Validity, and Limitations}
\label{sec:discussion}

The empirical findings from our primary frozen benchmark ($H_1$--$H_4$) and cross-schema study demonstrate that learned adaptive schema resolution with reactive recovery provides an efficient, reliable foundation for enterprise tool agents. Below, we analyze architectural boundaries, validity threats, and operational trade-offs.

\subsection{Schema Topologies and Structural Expressiveness}
Enterprise APIs exhibit diverse topologies, from flat records to deep hierarchies; our evaluation incorporated nested objects and arrays (e.g., line items in \texttt{create\_purchase\_order}, packages in \texttt{shipment\_creation}). Production MCP servers also employ polymorphic operators (\texttt{oneOf}, \texttt{anyOf}, \texttt{\$defs}). While JSON Pointer traversal handles nested definitions deterministically, mutually exclusive branches introduce pruning risks: pruning a discriminator field could misdirect validators. In LASR, discriminator fields and structural keys are designated as invariants ($C_I^*$), ensuring AST branching remains unambiguous. Automated DAG analysis for recursive schemas remains future work.

\subsection{Schema Validator Standards vs. Business Logic}
A core design philosophy of LASR-MCP is strict adherence to JSON Schema Draft 2020-12 specifications~\cite{jsonschema202012}. Standard validators extract RFC 6901 JSON Pointer paths~\cite{rfc6901} and keywords (\texttt{enum}, \texttt{pattern}, \texttt{minimum}, \texttt{minItems}) without bespoke heuristics. Enterprise backends also enforce business rules outside schema syntax (e.g., live inventory, credit checks). LASR-MCP explicitly scopes itself to \textit{syntactic and structural schema conformance}; transactional business failures must be handled by application compensation, not schema recomposition.

\subsection{Model-Family Sensitivities and Parameter Scale}
Controlled experiments benchmarked Google's \texttt{gemini-3.5-flash-lite} as a representative enterprise model. Pilot ablations observed unconstrained enum generation across evaluated configurations, showing the failure mode is not model-specific. For smaller open-weight models, lean schemas ($C_I^*$) prevent attention dilution across long contexts. For larger frontier models, transmitting monolithic schemas across multi-turn loops incurs compounding financial and latency costs. In both regimes, LASR’s lean baseline and closed-loop fallback provide operational advantages.

\subsection{Epistemological Separation and Benchmark Provenance}
We emphasize strict separation of empirical scopes: \textbf{1) Primary Frozen Benchmark ($N=32$ Tasks)}: Evaluates $H_1$--$H_4$ on the frozen procurement benchmark (commit \texttt{9807340}, SHA-256 \texttt{78c2c88e}), establishing predictive generalization for template transfer ($H_1$), invariant safety ($H_2$), reactive recovery lift ($H_3$), and workload token reduction ($H_4$). \textbf{2) Post-Freeze Robustness Study ($N=32$ Tasks)}: Evaluates cross-schema architecture transfer across three new synthetic schemas plus replication on the primary schema (four schemas total, SHA-256 \texttt{b3d57636}), providing preliminary evidence of mechanism robustness without altering the primary benchmark.

\section{Conclusion}
\label{sec:conclusion}
We introduced LASR-MCP, a learned adaptive schema resolution framework with dynamic in-tree recovery for enterprise tool-calling agents. By establishing an Invariant Foundation Floor ($C_I^*$), predicting task-conditioned constraint sets, and intercepting multi-error validator paths for native in-tree AST recomposition, LASR-MCP achieves a 22.1\% net workload-token reduction while attaining 100.0\% final structural validity on the primary frozen benchmark. Across four enterprise-style domains, LASR demonstrates 20.5\%--25.2\% token savings and complete 16/16 ambiguous-task recovery across evaluated cases, supporting the feasibility of a standards-compliant approach for cost-effective, reliable enterprise tool execution.


